\documentclass{article}

\usepackage{amsfonts}
\usepackage{amsthm}
\usepackage{mathrsfs}
\usepackage{amsmath}
\usepackage{graphicx}
\usepackage{hyperref}
\usepackage{stackengine}
\usepackage[mathscr]{eucal}
\newtheorem{theorem}{Theorem}
\theoremstyle{definition}\newtheorem{definition}{DEF}
\theoremstyle{lemma}\newtheorem{lemma}{Lemma}
\theoremstyle{corrolary}\newtheorem{corrolary}{Cor}
\setlength{\parindent}{0pt}

\begin{document}

\begin{definition}
Say the characteristic function of a random variable $X$ is
\begin{equation}
\varphi_X (t)=\mathbb{E}(e^{itX})
\end{equation}
\end{definition}

\begin{lemma}
If $X$ has an MGF, the characteristic function is
\begin{equation}
\varphi_X (t) = M_X (it)
\end{equation}
\end{lemma}
\begin{proof}
$$\varphi_X (t) = \mathbb{E}(e^{itX})=M_X(it)$$
\end{proof}

\begin{lemma}
Let $X$ be a random variable with PDF $f$.
Then the characteristic function is the Fourier transform
\begin{equation}
\varphi_X (t)=\mathscr{F}\{f\}\left(\frac{t}{2\pi}\right)
\end{equation}
\end{lemma}
\begin{proof}
By definition of the Fourier transform,
\begin{align*}
\mathscr{F}\{f\}\left(\frac{t}{2\pi}\right) &= \int_\mathbb{R} e^{\frac{2\pi itx}{2\pi}} f(x) dx\\
&= \int_\mathbb{R} e^{itx}f(x)dx\\
&=\mathbb{E}(e^{itX})\\
&=\varphi_X (t)
\end{align*}
\end{proof}

\begin{theorem}
Every random variable has a characteristic function.
\end{theorem}
\begin{proof}
Let $X$ be a random variable.
By Euler's formula, $|e^{itx}|=1$.
So,
\begin{align*}
\mathbb{E}(|e^{itX}|) &= \mathbb{E}(1)\\
&= 1\\
&<\infty
\end{align*}
Thus, by the definition of the Lebesgue integral, the characteristic function exists.
\end{proof}
\begin{corrolary}
$\varphi_X (t)=\varphi_Y (t)$ if and only if $X$ and $Y$ have the same distribution.
\end{corrolary}

\begin{theorem}
$\varphi_X(0)=1$ and $|\varphi_X(t)|\leq 1$.
\end{theorem}
\begin{proof}
\begin{align*}
\varphi_X(0)&=\mathbb{E}(e^{i(0)X})\\
&=\mathbb{E}(1)\\
&=1
\end{align*}
Also, absolute value is a convex function.
So by Jensen's inequality,
\begin{align*}
|\varphi_X (t)|&=|\mathbb{E}(e^{itX})|\\
&\leq \mathbb{E}(|e^{itX}|)\\
&=\mathbb{E}(1)\\
&=1
\end{align*}
\end{proof}

\begin{theorem}
Characteristic functions are uniformly continuous.
\end{theorem}
\begin{proof}
First, notice that for any $x,y\in \mathbb{R}$,
\begin{align*}
|e^{ix}-e^{iy}|&=|\int_x^y e^{it}dt|\\
&\leq \int_x^y |e^{it}| dt\\
&= y-x
\end{align*}
Assuming $\mathbb{E}(X)$ exists, we can use this inequality to see $e^{ixh}-1|\leq |hx|$.
Thus,
\begin{align*}
|\varphi(t+h)-\varphi(t)|&=|\mathbb{E}(e^{iX(t+h)}-e^{itAX})|\\
&\leq \mathbb{E}(|e^{itX}||e^{iXh}-1|)\\
&\leq \mathbb{E}(|hX|)\\
&=|h|\mathbb{E}(|X|)
\end{align*}
Conversely, if the expected value does not exist, $|e^{ith}-1|\rightarrow 0$ as $h\rightarrow 0$.
Thus, in either case, we have uniform continuity.
\end{proof}

\begin{lemma}
Let $X\sim \mathcal{N}(0, \sigma^2)$ and $Y\sim \mathcal{N}(0, \frac{1}{\sigma^2})$.
Then the PDF of $Y$ is $\frac{\sigma}{\sqrt{2\pi}} \varphi_X (t)$.
\end{lemma}
\begin{proof}
By properites of the characteristic function,
\begin{align*}
\varphi_X(t) &= M(it)\\
&=e^{\sigma^2 (it)^2 /2}\\
&= e^{-\sigma^2 t^2 /2}
\end{align*}
Also, the PDF of $Y$ is
\begin{equation*}
f_Y(t)=\frac{\sigma}{\sqrt{2\pi}} e^{-\sigma^2 t^2 /2}
\end{equation*}
Thus,
\begin{equation*}
f_Y(t)=\frac{\sigma}{\sqrt{2\pi}} \varphi_X(t)
\end{equation*}
\end{proof}

\textbf{Note-}
This lemma is related to measures of uncertainty in physics, including the Heisenberg's Uncertainty Principle.

\begin{definition}
Let $X$ be a random variable with characteristic function $\varphi_X$.
Say the $n^{th}$ cumulant of $X$ are the Taylor coefficients $K_n(X)$ such that
\begin{equation}
\ln \varphi_X(t)=\sum_{n=0}^\infty \frac{K_n(X) t^n}{n!}
\end{equation}
\end{definition}

\begin{theorem}
Let $X$ and $Y$ be independent random variables.
Then
\begin{equation}
K_n(X+Y) = K_N(X) + K_N(Y)
\end{equation}
\end{theorem}
\begin{proof}
Because $X$ and $Y$ are independent, $\varphi_{X+Y}(t)=\varphi_X(t)\varphi_Y(t)$.
Thus,
\begin{align*}
\ln \varphi_{X+Y}(t) &= \ln \varphi_X(t) \varphi_Y(t)\\
&= \ln \varphi_X(t) + \ln \varphi_Y(t)\\
&=\sum_{n=1}^\infty \frac{K_n(X)t^n}{n!} + \sum_{n=0}^\infty \frac{K_n(Y)t^n}{n!}\\
&=\sum_{n=0}^\infty \frac{[K_n(X)+K_n(Y)]t^n}{n!}
\end{align*}
\end{proof}

\end{document}
